
The copyright page looks great. The only thing missing is the printline, which should be set below the CIP data with space above the printline, i.e.,



at https://lccn.loc.gov...



10 9 8 7 6 5 4 3 2 1



The numbers should be separated with em-spaces.



Thanks,

Liz



Dear Tom,



Here follows another take on a possible promotional text for RMT4. The challenge is to provide useful information about the new edition without getting bogged down in unnecessary detail; while hopefully creating excitement in both new and old audiences.



Let�s see what you think.



Lars



Recursive methods provide powerful ways to pose and solve  problems in dynamic macroeconomics. Recursive Macroeconomic Theory provides both an introduction to recursive methods and more advanced material. Only practice in solving diverse problems fully conveys the advantages of the recursive approach, so this book offers  many applications. This fourth edition features two new chapters and substantial revisions to other chapters that
 demonstrate the power of recursive methods.



One new chapter applies the recursive approach to Ramsey taxation  and  sharply
 characterizes the  time inconsistency of optimal policies. These insights are used in other chapters to simplify recursive formulations of Ramsey  plans and credible government policies. The second new chapter explores the mechanics of matching models and identifies a  common channel though which productivity shocks operate across  a variety of matching models. Other chapters have been extended and refined. For example, there is new material on heterogenous beliefs both in complete and incomplete markets models; and  there is a deeper account of  forces that shape   aggregate labor supply elasticities in lifecycle models.



The book is suitable for both first- and second-year graduate courses in macroeconomics. Most chapters conclude with exercises; many exercises and examples use Matlab or Python computer programs.



